\documentclass[a4paper,12pt]{article}
\usepackage[utf8]{inputenc}
\usepackage[T1]{fontenc}
\usepackage[magyar]{babel}
\usepackage{amsmath, amssymb, amsthm}
\usepackage{geometry}
\usepackage{xcolor}
\usepackage{tcolorbox}
\usepackage{hyperref}
\usepackage{tikz}
\usetikzlibrary{shapes,arrows,positioning}

\geometry{margin=2.5cm}

\title{Téma Összefoglaló}
\author{AI Asszisztens}
\date{\today}

\begin{document}
\maketitle

\section{Elméleti áttekintés}

\begin{tcolorbox}[colback=blue!5!white,colframe=blue!75!black,title=Definíció]
Ide jön a definíció szövege...
\end{tcolorbox}

\begin{tcolorbox}[colback=green!5!white,colframe=green!75!black,title=Tétel]
Ide jön a tétel állítása...
\end{tcolorbox}

\begin{proof}
A tétel bizonyítása...
\end{proof}

\section{Kidolgozott példák}

\subsection{1. Példa}
Feladat leírása...

\medskip
\textbf{Megoldás:}
Lépésről lépésre...

\section{Önálló feladatok}
\begin{enumerate}
    \item Első feladat...
    \item Második feladat...
\end{enumerate}

\end{document}
